\SourceLeaf{314}%
by the modulus of Briggs's logarithms, with \(\rho\) taken in the sense specified in Article 16 (and hence \(\mu\) as well). For these constant factors we have
\[
\begin{aligned}
\log\mu&=7.9297527989\,(-20)\\
\log\nu&=4.9206912908\,(-30)
\end{aligned}
\]

\indent Through the second-order terms, these results agree with those given in Article 17. The sole purpose of the further expansion above was to make it clear that the second-order terms are entirely sufficient even for the most precise calculation. In fact, in the whole Hanoverian triangulation network there is no case in which the fourth-order terms reach two units in the tenth decimal place, and only a couple of cases in which they exceed \textit{one} unit in the tenth decimal place.

\Article{20}
If our formulas, which start from the mean latitude and azimuth at the two points instead of their values at one point, are to be used to solve the problem stated at the beginning of this memoir, the solution must proceed indirectly, or more precisely by successive approximation carried as far as desired. The procedure starts from any approximate value of \(T\) (if no other information or estimate is available, the given azimuth at the first point can serve as the initial value). A much more accurate value is derived from it, and the same calculation is repeated with that value. This continues until the results stabilize. In the first calculations, only 4 or 5 digits of the logarithms need to be taken into account, and \(\beta\) and \(\tau\) may be used in place of the corrected \(b\) and \(t\). Thus \(\lambda\) and \(l\) need not be considered yet in these first calculations. If the latitude at the first point is denoted by \(B^0\) and its azimuth by \(T^0\), the formulas are, in order:
\[
\begin{aligned}
\beta&=r\cos T\\
B&=B^0-\tfrac12\beta\\
\tau&=r\sin T\operatorname{tang}B\\
T&=T^0-\tfrac12\tau
\end{aligned}
\]
